\section{Ejemplos de contenido de Jordan}

\begin{eg}
  El conjunto $X = \Q \cap [a, b]$, con $a < b$, no tiene contenido nulo.
\end{eg}

\begin{proof}
  Como $\cl{X} = [a, b]$, para toda partición $P$ de un intervalo $I \supseteq [a, b]$ los subintervalos de $P$ que intersecan a $\cl{X}$ cubren $[a, b]$, y sus medidas suman al menos $b - a$. Luego $\overline{c}(X) \geq b - a > 0$.
\end{proof}

\begin{lemma}
  Sea $X \subseteq \R$ acotado. Si para todo $\e > 0$ existen un conjunto finito $F \subseteq X$ e intervalos abiertos $I_1, \dots, I_k$ tales que $X \setminus F \subseteq I_1 \cup \cdots \cup I_k$ y $|I_1| + \cdots + |I_k| < \e$, entonces $c(X) = 0$.
\end{lemma}

\begin{proof}
  Dado $\e > 0$, tomamos $F, I_1, \dots, I_k$ como en el enunciado con $\sum_{i=1}^k |I_i| < \frac{\e}{2}$. Si $F = \{x_1, \dots, x_q\}$, podemos cubrir $F$ con intervalos abiertos $I_{k+1}, \dots, I_n$ de longitud menor que $\frac{\e}{2q}$ cada uno. Entonces $X \subseteq I_1 \cup \cdots \cup I_n$ y $\sum_{i=1}^n |I_i| < \e$.

  Sean $\bar{I}_i = [\alpha_i, \beta_i]$ y $[a, b]$ un intervalo que contiene a $X$ y a todos los $\bar{I}_i$. Para $\eta > 0$, sea $P$ la partición de $[a, b]$ cuyos puntos son $a$, $b$ y todos los $\alpha_i$, $\beta_i$, $\alpha_i - \eta$, $\beta_i + \eta$ que están en $[a, b]$. Un subintervalo $R$ de $P$ que interseca a $\cl{X} \subseteq \bigcup_i \bar{I}_i$ interseca algún $\bar{I}_i$. Como $\alpha_i$ y $\beta_i$ son puntos de $P$, o bien $R \subseteq \bar{I}_i$, o bien $R$ comparte con $\bar{I}_i$ un único extremo, y en ese caso $R = [\beta_i, \beta_i']$ o $R = [\alpha_i', \alpha_i]$ con longitud a lo sumo $\eta$. Hay a lo sumo $2n$ subintervalos del segundo tipo, y los del primer tipo suman a lo sumo $\sum_i |I_i|$, pues sus interiores son disjuntos. Luego
  \[
    \overline{c}(X) \leq \overline{J}(X, P) \leq \sum_{i=1}^n |I_i| + 2n\eta < \e + 2n\eta.
  \]
  Como $\eta$ es arbitrario, $\overline{c}(X) \leq \e$, y como $\e$ es arbitrario, $\overline{c}(X) = 0$. Entonces $c(X) = 0$.
\end{proof}

\begin{eg}
  El conjunto de Cantor $C$ tiene contenido nulo.
\end{eg}

\begin{proof}
  Sea $F_n$ el conjunto obtenido en el paso $n$ de la construcción de $C$, unión de $2^n$ intervalos cerrados disjuntos de longitud $3^{-n}$. Entonces $C \subseteq F_n$ y la suma de las longitudes de los intervalos de $F_n$ es $1 - \frac{1}{3} \sum_{i=0}^{n-1} {\left( \frac{2}{3} \right)}^i = {\left( \frac{2}{3} \right)}^n$. Cada intervalo cerrado de $F_n$ está contenido en un intervalo abierto de longitud $2 \cdot 3^{-n}$, y la suma de las longitudes de éstos es $2 {\left( \frac{2}{3} \right)}^n$, que es menor que $\e$ para $n$ suficientemente grande. Por el lema (con $F = \varnothing$), $c(C) = 0$.
\end{proof}

\begin{theorem}
  Sea $S \subset \R^n$ acotado. Entonces $\overline{c}(\partial S) = \overline{c}(S) - \underline{c}(S)$. En particular, $S$ es medible Jordan si y sólo si $\partial S$ tiene contenido de Jordan nulo.
\end{theorem}

\begin{proof}
  Sea $I \subset \R^n$ un intervalo compacto que contiene a $\cl{S}$. Veamos primero que para toda partición $P$ de $I$,
  \[
    \overline{J}(\partial S, P) = \overline{J}(S, P) - \underline{J}(S, P).
  \]
  Un subintervalo $R$ de $P$ está contenido en $\interior S$ si y sólo si interseca a $\cl{S}$ y no interseca a $\partial S$. En efecto, si $R \subseteq \interior S$ entonces $R \cap \partial S = \varnothing$ y $R \cap \cl{S} \neq \varnothing$. Recíprocamente, si $R \cap \partial S = \varnothing$, entonces $R \subseteq \interior S \cup (\R^n \setminus \cl{S})$, unión de dos abiertos disjuntos; por ser $R$ conexo está contenido en uno de ellos, y si interseca a $\cl{S}$, está contenido en $\interior S$. Luego los subintervalos que intersecan a $\partial S = \cl{\partial S}$ son exactamente los que intersecan a $\cl{S}$ y no están contenidos en $\interior S$, lo que da la igualdad.

  Como $\overline{J}(S, P) \geq \overline{c}(S)$ y $\underline{J}(S, P) \leq \underline{c}(S)$, resulta $\overline{J}(\partial S, P) \geq \overline{c}(S) - \underline{c}(S)$ para toda $P$, y entonces $\overline{c}(\partial S) \geq \overline{c}(S) - \underline{c}(S)$.

  Para la desigualdad opuesta, dado $\e > 0$ sea $P_1$ tal que $\overline{J}(S, P_1) < \overline{c}(S) + \frac{\e}{2}$ y $P_2$ tal que $\underline{J}(S, P_2) > \underline{c}(S) - \frac{\e}{2}$. Sea $P = P_1 \cup P_2$ (refinamiento común). Como refinar achica las sumas exteriores y agranda las interiores,
  \begin{align*}
    \overline{c}(\partial S) \leq \overline{J}(\partial S, P) & = \overline{J}(S, P) - \underline{J}(S, P) \leq \overline{J}(S, P_1) - \underline{J}(S, P_2) \\
                                                              & < \overline{c}(S) - \underline{c}(S) + \e.
  \end{align*}
  Como $\e > 0$ es arbitrario, $\overline{c}(\partial S) \leq \overline{c}(S) - \underline{c}(S)$.

  Finalmente, $S$ es medible Jordan si y sólo si $\overline{c}(S) = \underline{c}(S)$, es decir, si y sólo si $\overline{c}(\partial S) = 0$, que equivale a $c(\partial S) = 0$ pues $0 \leq \underline{c}(\partial S) \leq \overline{c}(\partial S)$.
\end{proof}

\section{Medida de Lebesgue nula}

\begin{definition}[Medida de Lebesgue nula]
  Un conjunto $S \subseteq \R^n$ tiene \emph{medida de Lebesgue nula} si para todo $\e > 0$ puede cubrirse con una cantidad numerable de intervalos cuyas medidas suman menos que $\e$. Notamos $|S| = 0$.
\end{definition}

\begin{eg}
  Si $c(X) = 0$, entonces $|X| = 0$, pues existe una partición $P$ tal que los subintervalos que intersecan a $\cl{X}$ cubren a $X$ y suman menos que $\e$. Por ejemplo, el conjunto de Cantor tiene medida nula.
\end{eg}

\begin{prop}
  Toda unión numerable de conjuntos de medida nula tiene medida nula. En particular, todo conjunto numerable tiene medida nula; por ejemplo, $|\Q| = 0$.
\end{prop}

\begin{proof}
  Sean $F_1, F_2, \dots$ conjuntos de medida nula y $\e > 0$. Para cada $k \in \N$ cubrimos $F_k$ con una cantidad numerable de intervalos cuyas medidas suman menos que $\frac{\e}{2^k}$. La unión de todos estos intervalos es una familia numerable que cubre $\bigcup_k F_k$, y la suma de sus medidas es menor que $\sum_k \frac{\e}{2^k} = \e$.

  Todo conjunto con un único punto tiene medida nula (se cubre con un intervalo de medida menor que $\e$), de modo que todo conjunto numerable, por ejemplo $\Q = \{r_1, r_2, \dots\}$, tiene medida nula. Explícitamente, cubriendo a $r_k$ con el intervalo $\left( r_k - \frac{\e}{2^{k+2}}, r_k + \frac{\e}{2^{k+2}} \right)$ se obtiene una suma total $\sum_k \frac{\e}{2^{k+1}} = \frac{\e}{2} < \e$.
\end{proof}

TODO: terminar esta clase.

\section{Integrabilidad y medida de Jordan}

\section{Integración en conjuntos}
